\chapter{Reduced-Form Credit}\label{ch:rc:reduced-form-credit}

On 8 April 2009 the North American credit default swap market changed its contract
overnight. Until then every trade carried its own running spread, negotiated at
inception, and two swaps on the same name with the same maturity could pay different
coupons forever. From that day contracts on investment-grade names paid a fixed 100
basis points a year and those on high-yield names 500, and the difference between the
fixed coupon and the market's spread was settled in cash at the start. Converting a
quoted spread into that cash required a model everyone agreed on, and the industry
published one as open source. The model is the \omterm{def:rc:reduced-form-credit:reduced}{reduced-form model} of this chapter:
default is the first jump of a process with an intensity, the hazard rate of One
Quant Book 2, chapter 23. Book 2 priced a default swap with one flat hazard rate; here
the hazard rate has a term structure, the discounting comes from a real curve, and the
risks of a position, to the spread and to default itself, are measured.

\section{Default as the first jump of an intensity}

\begin{definition}[Reduced-form model]\label{def:rc:reduced-form-credit:reduced}
A \emph{reduced-form model}\index{reduced-form model} of credit makes the default time
$\tau_C$ of a name the first jump of a counting process with intensity $\lambda_t$ (a
Poisson process with a time-dependent or random rate; a Cox process of One Quant Book
4, chapter 7, when the intensity is random), without modelling why the firm defaults:
$\P(\tau_C\in[t,t+dt)\mid\tau_C>t) = \lambda_t\,dt$ and $Q_C(t) = \E[e^{-\int_0^t\lambda_s\,ds}]$.
\end{definition}

\begin{definition}[Probability of default, loss given default]\label{def:rc:reduced-form-credit:pd}
The \emph{probability of default}\index{probability of default} (PD) of a name over a
horizon $T$ is $1-Q_C(T)$. The \emph{loss given default}\index{loss given default}
(LGD) is the share of an exposure lost if it defaults, $\mathrm{LGD} = 1-R$, with $R$ the
recovery rate.
\end{definition}

\begin{definition}[Risky discount factor, risky annuity]\label{def:rc:reduced-form-credit:risky}
The \emph{risky discount factor}\index{risky discount factor} of a name to $T$ is
$P(0,T)Q_C(T)$, the value of one unit paid at $T$ if the name has survived, with
deterministic hazard rates and independent of rates. The \emph{risky
annuity}\index{risky annuity} of a default swap is the value of one unit a year of
premium paid while the name survives, $\sum_i\delta_iP(0,t_i)Q_C(t_i)$ plus the premium
accrued to a default within each period.
\end{definition}

\begin{proposition}[The credit triangle]\label{prop:rc:reduced-form-credit:triangle}
With a flat hazard rate $\lambda$, continuous premium payment and recovery $R$ paid at
default, the par spread of a default swap of any maturity is
$\mathcal S = \lambda(1-R)$.
\end{proposition}
\begin{proof}
Premium leg $\mathcal S\int_0^Te^{-(r+\lambda)t}dt$; protection leg
$(1-R)\int_0^T\lambda e^{-(r+\lambda)t}dt$. The integrals are equal; the par spread equates
the legs.
\end{proof}

\begin{definition}[Credit triangle]\label{def:rc:reduced-form-credit:triangledef}
The \emph{credit triangle}\index{credit triangle} is the rule of thumb $\mathcal S \approx
\lambda\times\mathrm{LGD}$ of \cref{prop:rc:reduced-form-credit:triangle}, used to read a
hazard rate from a spread (and back) in the head.
\end{definition}

\section{Pricing the legs of a default swap}

With quarterly premium dates $t_k$ and a default in $(t_{k-1},t_k]$ settled at $t_k$ with
half a period of accrued premium, the two legs on a \omterm{def:rc:multi-curve-and-collateral-discounting:curves}{discount curve} $P$ and a survival
curve $Q_C$ are
\begin{align*}
\text{premium} &= \mathcal S\sum_k\delta_kP(0,t_k)\Bigl(Q_C(t_k)+\tfrac12\bigl(Q_C(t_{k-1})-Q_C(t_k)\bigr)\Bigr),\\
\text{protection} &= (1-R)\sum_kP(0,t_k)\bigl(Q_C(t_{k-1})-Q_C(t_k)\bigr).
\end{align*}

The half-period accrual stands for the premium the buyer owes from the last payment date
to the default; the settlement at the period end instead of at default is a small
discounting error that the standard model removes by integrating over default times.
\Cref{fig:rc:reduced-form-credit:legs} shows the two legs as cash-flow streams.

\begin{omfigure}
\begin{tikzpicture}[x=1.05cm,y=1cm,font=\small]
\draw[->,thick] (0,0) -- (10.6,0) node[right] {$t$};
\foreach \k in {1,...,10} {\draw (\k,0.08) -- (\k,-0.08);}
\node[below] at (0,-0.1) {$0$};
\node[below] at (10,-0.1) {$T$};
\foreach \k in {1,...,6} {\draw[->,omDef,thick] (\k,0) -- (\k,-0.9);}
\node[omDef,align=left,anchor=west] at (0.1,-1.35) {premium $\mathcal S\,\delta$ each quarter while the name survives};
\draw[omThm,thick,dashed] (6.45,0.9) -- (6.45,-0.2);
\node[omThm,anchor=east] at (6.35,0.7) {default $\tau_C$};
\draw[->,omThm,very thick] (7,0) -- (7,1.6);
\node[omThm,align=left,anchor=west] at (7.15,1.3) {protection $1-R$ paid,\\ accrued premium owed};
\draw[omExo,dotted,thick] (7.2,-0.9) -- (10,-0.9);
\node[omExo,anchor=west] at (7.25,-0.6) {no further premium};
\end{tikzpicture}

\omcaption{The two legs of a default swap. The buyer pays the premium quarterly (blue)
until default; at default the seller pays the \omterm{def:rc:reduced-form-credit:pd}{loss given default} and the buyer the premium
accrued since the last payment (red). Schematic.}\label{fig:rc:reduced-form-credit:legs}
\end{omfigure}

\omcode{code/firm/cdscurve/firm_cdscurve.py}{40}{48}{The risky annuity and the protection leg on a discount curve and a hazard curve.}\label{lst:rc:reduced-form-credit:legs}

\section{Bootstrapping a hazard curve}

\begin{method}[Hazard curve from par spreads]\label{met:rc:reduced-form-credit:boot}
Take par spreads for increasing maturities $T_1<T_2<\dots$ and a recovery rate. Solve for
a constant hazard on $(0,T_1]$ that reprices the first; holding it, a constant hazard on
$(T_1,T_2]$ that reprices the second; and so on. The survival curve is
$Q_C(t) = e^{-\int_0^t\lambda}$ with the piecewise-constant hazards. Check that each input
reprices and that no hazard is negative (a negative hazard means the quotes are
inconsistent).
\end{method}

\omcode{code/firm/cdscurve/firm_cdscurve.py}{56}{68}{Bootstrapping piecewise-constant hazards so that every par default swap reprices.}\label{lst:rc:reduced-form-credit:boot}

\begin{example}[An investment-grade curve]\label{ex:rc:reduced-form-credit:curve}
Par spreads of 60, 75, 90, 120, 135 and 150 basis points at one, two, three, five, seven
and ten years (illustrative), recovery 40\%, on chapter~1's SOFR curve, bootstrap to
hazards of 1.00\%, 1.51\%, 2.04\%, 2.86\%, 3.03\% and 3.32\% on the successive intervals
(\cref{fig:rc:reduced-form-credit:hazard}). The \omterm{def:rc:reduced-form-credit:triangledef}{credit triangle} gives 1.00\%, 1.25\%,
1.50\%, 2.00\%, 2.25\% and 2.50\%: the average hazard to each maturity, where the
bootstrap gives the forward hazard on each interval. The \omterm{def:rc:reduced-form-credit:pd}{probability of default} is
0.995\% over one year, 9.76\% over five and 23.1\% over ten
(\cref{fig:rc:reduced-form-credit:survival}).
\end{example}

\begin{omfigure}
\begin{tikzpicture}
\begin{axis}[width=11.5cm,height=5.2cm,
  xlabel={time (years)},ylabel={hazard rate (\%)},
  tick label style={font=\small},label style={font=\small},
  xmin=0,xmax=10,ymin=0,ymax=3.8,grid=major,grid style={gray!25},
  legend columns=2,legend style={font=\footnotesize,at={(0.5,-0.25)},anchor=north,draw=none,fill=none,/tikz/every even column/.append style={column sep=8pt}},
  table/col sep=comma]
\addplot[omThm,very thick,const plot] table[x=t,y=hazard]{figdata/rates-credit-risk/13-reduced-form-credit/hazard.csv};
\addplot[omDef,thick,dashed,const plot] table[x=t,y=triangle]{figdata/rates-credit-risk/13-reduced-form-credit/hazard.csv};
\legend{bootstrapped (forward) hazard,credit triangle $\mathcal S/(1-R)$ at each maturity}
\end{axis}
\end{tikzpicture}

\omcaption{The hazard curve bootstrapped from six par spreads, and the average hazard
the \omterm{def:rc:reduced-form-credit:triangledef}{credit triangle} reads from each spread. An upward-sloping spread curve means higher
forward hazards further out. Data: illustrative spreads on chapter~1's curve; the
chapter's tutorial.}\label{fig:rc:reduced-form-credit:hazard}
\end{omfigure}

\begin{omfigure}
\begin{tikzpicture}
\begin{axis}[width=11.5cm,height=5.2cm,
  xlabel={horizon (years)},ylabel={probability (\%)},
  tick label style={font=\small},label style={font=\small},
  xmin=0,xmax=10,ymin=0,ymax=100,grid=major,grid style={gray!25},
  legend columns=2,legend style={font=\footnotesize,at={(0.5,-0.25)},anchor=north,draw=none,fill=none,/tikz/every even column/.append style={column sep=8pt}},
  table/col sep=comma]
\addplot[omDef,very thick] table[x=t,y=q]{figdata/rates-credit-risk/13-reduced-form-credit/survival.csv};
\addplot[omThm,very thick,dashed] table[x=t,y=pd]{figdata/rates-credit-risk/13-reduced-form-credit/survival.csv};
\legend{survival probability $Q_C(t)$,probability of default $1-Q_C(t)$}
\end{axis}
\end{tikzpicture}

\omcaption{Survival and default probabilities of the bootstrapped name to ten years. The
survival curve bends down faster as the forward hazard rises; the \omterm{def:rc:reduced-form-credit:pd}{probability of default}
reaches 9.8\% at five years and 23.1\% at ten. Data: the chapter's tutorial.}\label{fig:rc:reduced-form-credit:survival}
\end{omfigure}

\begin{remark}[Recovery is a convention, expected loss is not]\label{rem:rc:reduced-form-credit:recovery}
Spreads pin down the expected loss, roughly $\lambda(1-R)$, not the hazard alone. With a
recovery of 25\% instead of 40\% the same spreads bootstrap to hazards of 0.80\% to 2.64\%
and a five-year \omterm{def:rc:reduced-form-credit:pd}{probability of default} of 7.87\% instead of 9.76\%, while every quoted
swap is worth the same. Recovery matters for positions that are not par swaps: a bond
priced near par, a swap far from its coupon, and the jump-to-default of any of them.
\end{remark}

\begin{remark}[Inconsistent quotes]\label{rem:rc:reduced-form-credit:negative}
If the three-year spread were 50 basis points, below the two-year 75, repricing it would
need a negative hazard on the third year: the bisection of \cref{lst:rc:reduced-form-credit:boot}
stops at zero and reprices only 51.0 basis points. A curve that fails this check is
telling the desk that one quote is stale, or that the recovery assumed is wrong.
\end{remark}

\section{The standard model and its conventions}

\begin{definition}[ISDA standard model]\label{def:rc:reduced-form-credit:isda}
The \emph{ISDA standard model}\index{ISDA standard model} is the open-source
reduced-form pricer, with agreed conventions (a flat hazard implied from a quoted
spread, a standard recovery, the day's standard \omterm{def:rc:multi-curve-and-collateral-discounting:curves}{discount curve}, accrual and payment
rules), that the market uses to convert between a contract's quoted spread and the
upfront payment at its standard coupon.
\end{definition}

A quoted spread of a standard contract is thus not a par spread of a bootstrapped curve:
it is the flat spread that, in the standard model, gives the traded upfront. The two
agree for the maturity quoted and differ for positions of other maturities, which a
desk values on its bootstrapped curve.

\begin{example}[Upfront at a standard coupon]\label{ex:rc:reduced-form-credit:upfront}
The five-year contract quoted at 120 basis points with a 100-basis-point coupon costs
the protection buyer an upfront of 0.872\% of notional (flat hazard from the quote, the
five-year zero rate of 3.42\% flat, recovery 40\%). A name quoted at 500 has no upfront at
the 500 coupon and one of 15.06\% at the 100 coupon (\cref{fig:rc:reduced-form-credit:upfront}).
\end{example}

\begin{omfigure}
\begin{tikzpicture}
\begin{axis}[width=11.5cm,height=5.2cm,
  xlabel={quoted spread (bp)},ylabel={upfront (\% notional)},
  tick label style={font=\small},label style={font=\small},
  xmin=0,xmax=800,ymin=-25,ymax=25,grid=major,grid style={gray!25},
  legend columns=2,legend style={font=\footnotesize,at={(0.5,-0.25)},anchor=north,draw=none,fill=none,/tikz/every even column/.append style={column sep=8pt}},
  table/col sep=comma]
\addplot[omDef,very thick] table[x=s,y=u100]{figdata/rates-credit-risk/13-reduced-form-credit/upfront.csv};
\addplot[omThm,very thick,dashed] table[x=s,y=u500]{figdata/rates-credit-risk/13-reduced-form-credit/upfront.csv};
\legend{100 bp coupon,500 bp coupon}
\end{axis}
\end{tikzpicture}

\omcaption{Upfront paid by the protection buyer of a five-year standard contract, from
the quoted spread, at the two North American standard coupons. The upfront is zero
where the quote equals the coupon and negative (received) below it. Data: the standard
flat-hazard conversion; the chapter's tutorial.}\label{fig:rc:reduced-form-credit:upfront}
\end{omfigure}

\begin{dated}{2026-09}{Standard contracts}\label{dat:rc:reduced-form-credit:standard}
Since the Big Bang Protocol of 8 April 2009, North American single-name default swaps
trade with a 100-basis-point coupon on investment-grade names and 500 on high-yield
names, with an upfront settling the difference, and settle credit events by auction.
The standard model used for the conversion is published as open-source code under an
ISDA licence and administered by S\&P Global Market Intelligence.
\end{dated}

\section{Bonds on the hazard curve}

The same survival curve prices a fixed-coupon bond of the name: each coupon is a
\omterm{def:rc:reduced-form-credit:risky}{risky discount factor}, and default pays the recovery rate on the face value (recovery
of par) at the period end, as in \texttt{risky\_bond}. The value of a bond on the hazard
curve bootstrapped from default swaps, compared with its market price, gives the
difference between the two markets' prices of the same credit.

\begin{example}[A five-year bond on the curve]\label{ex:rc:reduced-form-credit:bond}
A five-year 5\% bond of the name (semi-annual coupons) is worth 107.08 per 100 on the
riskless curve and 101.38 on the bootstrapped hazard curve. The \omterm{def:rc:reduced-form-credit:risky}{risky discount factor} to
five years is 0.7604, against a riskless 0.8427. If the bond trades at 97.00, the flat
hazard that reprices it corresponds to a default-swap spread of 220 basis points: the
bond is 100 basis points cheaper than the default swap says, a negative basis of
$-99.8$ basis points (default swap minus bond-implied spread).
\end{example}

A negative basis invites buying the bond and buying protection: the package collects
the bond's excess spread while the protection covers default. It is not an arbitrage: it
must be funded, it is marked to a basis that can widen, and its default payoff depends
on the hedge ratio and on the recovery; the weekend problem measures it.

\section{Credit risk measures: CS01 and jump to default}

\begin{definition}[CS01]\label{def:rc:reduced-form-credit:cs01}
The \emph{CS01}\index{CS01} of a position is the change in its value when a credit
spread rises by one basis point: bucketed, per quoted maturity with the hazard curve
re-bootstrapped, or parallel.
\end{definition}

\begin{definition}[Jump-to-default risk]\label{def:rc:reduced-form-credit:jtd}
\emph{Jump-to-default risk}\index{jump-to-default risk} (JTD) is the change in a
position's value if the name defaults now: the protection paid or received,
$\pm(1-R)N$, minus the position's current value, which disappears. It is not captured by
\omterm{def:rc:reduced-form-credit:cs01}{CS01}: default is a jump, not a large spread move.
\end{definition}

\begin{example}[A protection position]\label{ex:rc:reduced-form-credit:position}
Long USD~10 million of five-year protection at the 100 coupon on the curve of
\cref{ex:rc:reduced-form-credit:curve}: worth USD~88\,133 (the \omterm{def:rc:reduced-form-credit:risky}{risky annuity} is 4.41),
with a \omterm{def:rc:reduced-form-credit:cs01}{CS01} of USD~4\,391 on the five-year quote and almost nothing elsewhere, and a
jump-to-default gain of USD~5.91 million.
\end{example}

\begin{definition}[Default risk premium]\label{def:rc:reduced-form-credit:premium}
The \emph{default risk premium}\index{default risk premium} is the difference between
the risk-neutral default probability implied by spreads and the real-world probability
estimated from default histories: spreads pay for expected losses, for their
uncertainty and for the illiquidity of credit, so implied PDs exceed historical ones,
most for high-quality names.
\end{definition}

\begin{remark}[Which PD to use]\label{rem:rc:reduced-form-credit:which}
Pricing and hedging use the risk-neutral PD (9.76\% over five years above); capital
models, provisions and ratings use real-world PDs. If the real-world five-year PD of the
name were 1.5\%, spreads would price more than six times the expected default losses. Hull,
Predescu and White (2005) found seven-year risk-neutral default intensities 16.8 times the historical
ones for Aaa bonds, 9.8 for A, 5.1 for Baa, 2.1 for Ba and 1.2 for B.
\end{remark}

\section{Tutorial: a hazard curve}\label{tut:rc:reduced-form-credit:tutorial}

\begin{tutorial}
\textbf{Goal.} Bootstrap a hazard curve from par spreads, convert quotes to upfronts, and
measure a position's \omterm{def:rc:reduced-form-credit:cs01}{CS01} and jump-to-default. \textbf{End state:}
\cref{fig:rc:reduced-form-credit:hazard,fig:rc:reduced-form-credit:upfront} and the
numbers of \cref{ex:rc:reduced-form-credit:curve,ex:rc:reduced-form-credit:position}.
\begin{tutsteps}
\item \textbf{Bootstrap}: \texttt{bootstrap(disc, tenors, spreads, recovery)}; check each
par spread reprices.
\item \textbf{Triangle}: \texttt{hazard\_table()} against $\mathcal S/(1-R)$.
\item \textbf{Upfronts}: \texttt{standard\_upfront} (Book~2's flat-hazard conversion).
\item \textbf{Risk}: \texttt{protection\_position()}; \texttt{fig\_rc\_credit.py} writes the
charts.
\end{tutsteps}
\textbf{What to change next.} Use a recovery of 25\% and compare hazards and upfronts;
make the three-year spread 50 basis points and watch the bootstrap fail to reprice it
(\cref{rem:rc:reduced-form-credit:negative}).
\end{tutorial}

\section{Build: the credit curve}\label{bld:rc:reduced-form-credit:cdscurve}

\begin{build}
\bfield{Purpose} The credit curves of the miniature firm: every default swap, credit
bond and counterparty in chapters 14 to 20 reads a hazard curve built here.
\bfield{Interface} \texttt{HazardCurve(times, hazards)} with \texttt{cum} and
\texttt{survival}; \texttt{legs}, \texttt{par\_spread}, \texttt{bootstrap};
\texttt{standard\_upfront}; \texttt{cds\_value}; \texttt{risky\_bond};
\texttt{bucketed\_cs01}; \texttt{jump\_to\_default}.
\bfield{Rules} Quarterly premiums with half-period accrual on default, defaults settled at
period ends; any \omterm{def:rc:multi-curve-and-collateral-discounting:curves}{discount curve} with \texttt{df\_t}; Book~2's \texttt{firm\_cds} for the
flat conversion.
\bfield{Acceptance tests} \texttt{code/firm/cdscurve/tests/}: every quote reprices; one
tenor gives Book~2's flat hazard; the upfront vanishes at the coupon; a bond with zero
hazard is a riskless bond; \omterm{def:rc:reduced-form-credit:cs01}{CS01} falls on the position's own maturity; jump-to-default of
a protection buyer is $(1-R)N$ minus its value.
\bfield{Stretch} Business-day schedules and the standard model's exact conventions
(step-in, protection from trade date); recovery-rate risk; stochastic intensity.
\end{build}

\begin{omsources}
\item Paul, Weiss, ``The Big Bang Protocol and a new structural framework for credit
default swaps'', 24 March 2009.
\item ISDA CDS Standard Model, \texttt{cdsmodel.com}.
\item D.\ O'Kane, \emph{Modelling Single-name and Multi-name Credit Derivatives}, Wiley,
2008, chapters 3--7.
\item J.\ Hull, M.\ Predescu and A.\ White, ``Bond prices, default probabilities and risk premiums'',
\emph{Journal of Credit Risk} 1(2), 2005, 53--60.
\end{omsources}

\section{Exercises}

\begin{exercise}[$\star$]\label{exo:rc:reduced-form-credit:1}
A name's five-year spread is 180 basis points and recovery 40\%. Estimate its hazard rate
and five-year survival probability with the \omterm{def:rc:reduced-form-credit:triangledef}{credit triangle}.
\end{exercise}

\begin{exercise}[$\star$]\label{exo:rc:reduced-form-credit:2}
Why is a bootstrapped hazard higher than the triangle's hazard when the spread curve
slopes upward?
\end{exercise}

\begin{exercise}[$\star$]\label{exo:rc:reduced-form-credit:3}
In \cref{ex:rc:reduced-form-credit:position}, what is the jump-to-default of the seller of
the same protection?
\end{exercise}

\begin{exercise}[$\star\star$]\label{exo:rc:reduced-form-credit:4}
Estimate the \omterm{def:rc:reduced-form-credit:cs01}{CS01} of the protection position from its \omterm{def:rc:reduced-form-credit:risky}{risky annuity}, and compare with
USD~4\,391.
\end{exercise}

\begin{exercise}[$\star\star$]\label{exo:rc:reduced-form-credit:5}
A high-yield name is quoted at 350 basis points. Give the five-year upfront at the 500
coupon with the chapter's conventions, say who pays it, and give the upfront at the 100
coupon.
\end{exercise}

\begin{exercise}[$\star\star$]\label{exo:rc:reduced-form-credit:6}
Why is jump-to-default not the \omterm{def:rc:reduced-form-credit:cs01}{CS01} times a very large spread move?
\end{exercise}

\begin{exercise}[$\star\star\star$]\label{exo:rc:reduced-form-credit:7}
\emph{Coding.} Compute the negative-basis package of the weekend problem: the
bond-implied spread, the basis, and the package's jump-to-default with a par-for-par
hedge.
\end{exercise}

\begin{exercise}[$\star\star\star$]\label{exo:rc:reduced-form-credit:8}
\emph{Find the flaw.} ``The five-year spread is 120 basis points, so the market thinks
the name has a 2\% chance of defaulting each year.''
\end{exercise}

\section{Problem: The Negative Basis Package}

\begin{problem}[{Weekend problem --- buying a cheap bond with protection}]\label{pb:rc:reduced-form-credit:1}
The chapter's name has a five-year 5\% bond (semi-annual) trading at 97.00 while its
five-year default swap is quoted at 120 basis points. A fund buys USD~100 million face of
the bond and five-year protection at the 100 coupon.

\medskip\noindent\textbf{Part I --- The basis.}
\begin{enumerate}
\item Give the bond's value on the bootstrapped hazard curve.
\item Give the flat par spread that reprices the bond at 97.00.
\item Give the basis (default swap minus bond-implied spread).
\item Why might a bond trade so far below its default-swap-implied value?
\item What does the fund earn if nothing happens?
\end{enumerate}

\medskip\noindent\textbf{Part II --- Hedge ratios.}
\begin{enumerate}[resume]
\item Give the jump-to-default with USD~100 million of protection (par for par).
\item Give it with USD~97 million (market value).
\item Give the package's \omterm{def:rc:reduced-form-credit:cs01}{CS01} in each case.
\item Give the recovery at which the market-value package breaks even on default.
\item Which hedge would you choose, and why?
\end{enumerate}

\medskip\noindent\textbf{Part III --- Risks.}
\begin{enumerate}[resume]
\item What happens to the package if the basis widens further?
\item What funding risk does the fund carry?
\item What is the cheapest-to-deliver option in the default swap worth to the fund?
\item What counterparty risk does the protection carry?
\item How would the package behave in a 2008-style dash for cash?
\end{enumerate}

\medskip\noindent\textbf{Part IV --- Judgement.}
\begin{enumerate}[resume]
\item Why is the negative basis called an arbitrage, and why is it not one?
\item Which desks run it, and with what balance sheet?
\item What would you ask before putting it on?
\item State the \emph{named result}: the basis, and the package's jump-to-default and \omterm{def:rc:reduced-form-credit:cs01}{CS01}
with the par-for-par hedge.
\item In one sentence: what does a negative basis measure?
\end{enumerate}
\end{problem}

\section{Interview questions}

\begin{interviewq}[$\star$ \iqroles{trader, researcher}]\label{iq:rc:reduced-form-credit:1}
State the \omterm{def:rc:reduced-form-credit:triangledef}{credit triangle} and use it on a 200-basis-point spread.
\end{interviewq}

\begin{interviewq}[$\star\star$ \iqroles{researcher, developer}]\label{iq:rc:reduced-form-credit:2}
How do you bootstrap a hazard curve, and what can go wrong?
\end{interviewq}

\begin{interviewq}[$\star\star$ \iqroles{trader, bank}]\label{iq:rc:reduced-form-credit:3}
What changed with the 2009 standardisation of default swaps, and why did it need a
common model?
\end{interviewq}

\begin{interviewq}[$\star\star$ \iqroles{risk}]\label{iq:rc:reduced-form-credit:4}
Why report jump-to-default separately from \omterm{def:rc:reduced-form-credit:cs01}{CS01}?
\end{interviewq}

\begin{interviewq}[$\star\star\star$ \iqroles{researcher, risk}]\label{iq:rc:reduced-form-credit:5}
Risk-neutral default probabilities are several times historical ones. Explain, and say
which you would use for pricing, provisioning and capital.
\end{interviewq}

\begin{interviewq}[$\star\star\star$ \iqroles{trader}]\label{iq:rc:reduced-form-credit:6}
A bond trades 100 basis points wide of its default swap. How do you trade it, and what
could go wrong?
\end{interviewq}
